\section*{2 Preliminaries}
\paragraph{Definition3 (Polynomial Calculus/Polynomial Calculus Resolution).}
Let $\Gamma=\{P_1\ldots P_m\}$ be an unsolvable system of polynomials in variables $\{x_1\ldots x_n\}$ over $\mathbb F$. A PC (Polynomial Calculus) refutation of $\Gamma$ is a sequence of polynomials $\{R_1\ldots R_s\}$ such that $R_s=1$ and for every $\ell\in[s]$, $R_{\ell}\in\Gamma$, $R_{\ell}$ is either a polynomial from $\Gamma$, or is obtained from two previous polynomials $R_j,R_k$, $j,k<\ell$ by one of the following derivation rules:
\[R_{\ell}=\alpha R_j+\beta R_k\text{ for }\alpha,\beta\in\mathbb F,\qquad R_{\ell}=x_iR_k\text{ for some }i\in[n].\]
The size of the refutation is $\sum_{\ell=1}^s|R_{\ell}|$, where $|R_{\ell}|$ is the number of monomials in the polynomial $R_{\ell}$. The degree of the refutation is $\max_{\ell}\deg(R_{\ell})$.
A PCR (Polynomial Calculus Resolution) refutation is a PC refutation over the set of Boolean variables $\{x_1\ldots x_n,\bar x_1\ldots\bar x_n\}$ where $\{\bar x_1\ldots\bar x_n\}$ are twin variables of $\{x_1\ldots x_n\}$ i.e. the equations $x_i^2-x_i=0$, $\bar x_i^2-\bar x_i=0$ and $x_i\bar x_i=0$ are treated as axioms.
\paragraph{Definition4 (PC plus Extension Axioms).}
Let $\Gamma=\{P_1\ldots P_m\}$ be a set of polynomials in variables $\{x_1\ldots x_n\}$ over a field $\mathbb F$. We will refer to the polynomials in $\Gamma$ as (initial) axioms. Let $\mathbf z=z_1\ldots z_M$ be new extension variables with corresponding extension axioms $z_j-Q_j(x_1\ldots x_n)$. A PC+Ext (PC plus extension) refutation of $\Gamma$ with $M$ extension axioms $\mathsf{Ext}=\{z_j-Q_j(x_1\ldots x_n)\}$ is a PC refutation of the set of polynomials $\Gamma'=\{P_1\ldots P_m,z_1-Q_1\ldots z_M-Q_M\}$. An extension axiom $z_j=Q_j(x_1\ldots x_n)$ is $\kappa$-local if $Q_j$ is a $\kappa$-junta; that is, if the polynomial $Q_j$ defining $z_j$ involves at most $\kappa$ of the $x$-variables. We say that $\Pi$ is a $(M,\kappa)$-PC+Ext refutation of $\Gamma$ if it is a PC+Ext refutation of $\Gamma$ with $M$ extension axioms, each of which are $\kappa$-local. The size of the refutation is total size of all lines in the refutation, including the polynomials in $\Gamma$ plus the extension axioms (where the size of a line $P\in\Pi$ is the number of monomials in $P$).
We note that our definition of extension axioms is more limited than the general notion of extension axioms. Here we only allow the extension variables to depend on the original variables from $\Gamma$; the more general definition allows the extension variables to depend on the original $x$-variables, and also on other extension variables.
\paragraph{Definition5 ($k$-local CSPs).}
A constraint $C_i$ over Boolean variables $\{x_1,\ldots,x_n\}$ is simply a Boolean formula over these variables. $C_i$ is a $k$-local constraint if $C_i$ depends on at most $k$ variables. A $k$-CSP $C=C_1\land\ldots\land C_m$ over $\{x_1,\ldots,x_N\}$ is the conjunction of a set of $k$-local constraints.
We translate a $k$-CSP formula into a system of polynomial equations using the standard PCR translation which we define next.
\paragraph{Definition6 (Converting $k$-CSPs into Polynomial Equations).}
Let $C$ be a $k$-local constraint over variables $x_{i_1},\ldots,x_{i_k}$. We convert $C$ to a polynomial equation, $p(C)$, using the translations $p(x_{i_j})=1-x_{i_j}$, $p(\neg A)=1-p(A)$, $p(A\lor B)=p(A)\cdot p(B)$. It is easy to check that for any Boolean assignment $\alpha$ to the underlying variables, $C(\alpha)=1\leftrightarrow p(\alpha)=0$, and $C(\alpha)=0\leftrightarrow p(\alpha)=1$.
A $k$-CSP $C=C_1\land\ldots\land C_m$ over $\{x_1,\ldots,x_n\}$ converts to a set of polynomial equations $\{E_j\mid j\in[m]\}\cup\{B_i\mid i\in[n]\}$ over $\{x_1,\ldots,x_n\}\cup\{\bar x_1,\ldots,\bar x_n\}$ where $E_j$ is the polynomial equation $p(C_j)$. In addition, we add the Boolean axioms $\{B_i\mid i\in[n]\}$, where $B_i=\{x_i^2-x_i=0,\bar x_i^2-\bar x_i=0,x_i\bar x_i=0\}$ which force $x_i,\bar x_i$ to be zero-one valued, and force exactly one of $x_i,\bar x_i$ to be one.
\section*{3 The Hard Formulas}
We distinguish between the case $p=2$ and the case $p>2$, and concentrate on the latter. This is because the case $p=2$ does not require any new technical ideas, and we can pick from a large number of known hard tautologies for this case, such as random CNF's. Over $\mathbb F_2$, every extension variable is zero-one valued, and so standard size-degree tradeoffs pertain even with respect to extension variables. Also, $\kappa$-local extension variables can change the degree by at most a factor of $\kappa$, therefore a degree lower bound of $\Omega(n)$ for the original tautology over $n$ variables implies a degree lower bound of $\Omega(n/\kappa)$ after adding $\kappa$-local extension variables. Known size-degree tradeoffs imply that the degree must be at least square root of the number of variables in order to obtain exponential size lower bounds, this immediately gives a lower bound tolerating close to $n^2/\kappa^2$ many $\kappa$-local extension variables [10,6,17].
Over any field, there are unsatisfiable families of $k$-CNF formulas (e.g. the Tseitin tautologies as well as random parity equations) that require linear degree but have polynomial sized proofs with a linear number of extension variables [8,6]. Therefore formulas that require high PC degree are not sufficient. Instead we will create our hard examples by taking a hard instance and then using selector variables to pick out a subset of the constraints. Similar ideas were used earlier (e.g., [12]). In more detail, our underlying hard unsatisfiable formulas, $\{F^{\mathsf{SEL}}_{n,k}\}$, will be constructed from a family of $k$-CSP formulas, $F_{n,k}$, that have the property that any sufficiently large subset of the constraints of $F_{n,k}$ is unsatisfiable and still requires large PC degree.
\paragraph{Definition7.}
Let $F_{n,k}=\{E_j\mid j\in[m]\}\cup\{B_i\mid i\in[n]\}$ be the system of degree-$k$ polynomial equations over $\mathbf x=\{x_i,\bar x_i\mid i\in[n]\}$, obtained by converting a size-$m$ $k$-CSP as given by Definition6. For convenience, we will index the polynomial equations $E_j$ in binary notation, so for example if $b_1\ldots b_{\log m}\in\{0,1\}^{\log m}$ is the binary notation for $j\in[m]$, we will write $E_j$ as $E_{b_1\ldots b_{\log m}}$. We define a new set of polynomial equations $F^{\mathsf{SEL}}_{n,k}$ with parameters $m,m'$ as follows. The variables are $\mathbf x\cup\mathbf y$, where $\mathbf x$ are the original variables of $F_{n,k}$ and $\mathbf y=\{y_{i,j},\bar y_{i,j}\mid i\in[m'],j\in[\log m]\}$ are new ``pigeon'' variables. Let $E^{\mathsf{SEL}}$ be the following set of equations, where $\mathbf y_i\neq b_1\ldots b_{\log m}$ abbreviates the monomial $\prod_{b_j=1}y_{i,j}\prod_{b_j=0}\bar y_{i,j}$:
\begin{enumerate}[label=(\roman*)]
\item $\forall i\in[m'],\ \forall b_1\ldots b_{\log m}\in\{0,1\}^{\log m},\ (\mathbf y_i\neq b_1\ldots b_{\log m})\cdot E_{b_1\ldots b_{\log m}}=0$;
\item $\forall i,i'\in[m'],i\neq i',\ \forall b_1\ldots b_{\log m}\in\{0,1\}^{\log m},$
\[(\mathbf y_i\neq b_1\ldots b_{\log m})\cdot(\mathbf y_{i'}\neq b_1\ldots b_{\log m})=0.\]
\end{enumerate}
$F^{\mathsf{SEL}}_{n,k}$ consists of the polynomial equations $E^{\mathsf{SEL}}$ together with the Boolean axioms $B_{i,j}=\{y_{i,j}^2-y_{i,j}=0,\bar y_{i,j}^2-\bar y_{i,j}=0,y_{i,j}\bar y_{i,j}=0\}$ for all $i\in[m'],j\in[m]$.
Intuitively we think of the $y$ variables as a mapping from $m'$ pigeons to $m$ holes, where the holes correspond to the $m$ axioms/constraints from $E$. For $i\in[m']$, the $i$th ``pigeon'' $\mathbf y_i$ selects a hole (an equation from $E$).
The first set of polynomial equations in $E^{\mathsf{SEL}}$ states that if pigeon $\mathbf y_i$ selects the equation $E_{b_1\ldots b_{\log m}}$, then this equation must be satisfied; the second set of equations in $E^{\mathsf{SEL}}$ states that the mapping is one-to-one and thus altogether the $y$ selector variables choose a subset $E'$ of exactly $m'$ equations from $E$. Thus $F^{\mathsf{SEL}}_{n,k}$ asserts that there exists a subset of $m'$ constraints of $F_{n,k}$ (chosen by the $y$-variables) that are satisfiable.
Throughout this paper, the $x$-variables are the variables that underly $F_{n,k}$; the $y$-variables are the selector/pigeon variables described above that choose a subset of $m'$ constraints from $F_{n,k}$, and the extension variables used in the PC+Ext refutation will be the $z$-variables.
Our hard instances will be $F^{\mathsf{SEL}}_{n,k}$, with $m=10n$, $m'=(1-\epsilon)m$, where $F_{n,k}$ is (the polynomial translation of) an unsatisfiable $k$-CSP formula with $m=10n$ $k$-local constraints over variables $\mathbf x=x_1\ldots x_n$, satisfying the follow property:
\paragraph{Property8.} Every subset of $(1-\epsilon)m'$ constraints is unsatisfiable and requires linear PC degree.
The following Theorem shows that for sufficiently large $n$, such formulas exist. Similar proofs have appeared in several papers (e.g., [5]) but we give a proof in the Appendix for completeness.
\paragraph{Theorem9.} Let $m=10n$. Then there exists constants $k>0$, $0<\epsilon<1$ such that for sufficiently large $n$, there exists $k$-CSP formulas $\{F_{n,k}\}$ with $m$ constraints such that Property8 holds with $m'=(1-\epsilon)m$.
